\documentclass[runningheads]{llncs}

\usepackage[T1]{fontenc}
\usepackage[utf8]{inputenc}
\usepackage{graphicx}
\usepackage{amssymb}
\usepackage{booktabs}
\usepackage{hyperref}
\usepackage{enumitem}

\renewcommand\UrlFont{\color{blue}\rmfamily}
\urlstyle{rm}

\begin{document}
\sloppy

\title{Navegação Autónoma na Formula Student: Uma Abordagem Baseada em Aprendizagem por Reforço}
\titlerunning{Navegação Autónoma na Formula Student com RL}

\author{Luís Miguel Pereira Silva\inst{1} \and
Pedro Miguel S.\ A.\ Urbano dos Reis\inst{1}}

\authorrunning{L. Silva e P. Reis}

\institute{Mestrado em Inteligência Artificial, Universidade do Minho, Braga, Portugal\\
\email{\{pg60390,pg59908\}@alunos.uminho.pt}}

\maketitle

\begin{abstract}
Este documento apresenta o planeamento de um trabalho prático centrado no desenvolvimento de um sistema de condução autónoma para a competição Formula Student, recorrendo a técnicas de Aprendizagem por Reforço (RL). O problema consiste em treinar um agente capaz de navegar numa pista delimitada por cones coloridos, respeitando as regras da categoria \textit{Driverless} (FS-AI). Será desenvolvido um ambiente de simulação 2D com modelo cinemático de bicicleta, geração procedural de pistas e penalizações realistas. A metodologia assenta na implementação de um agente \textit{Soft Actor-Critic} (SAC) para controlo em espaços de ação contínuos, utilizando \textit{Proximal Policy Optimization} (PPO) como \textit{baseline} comparativo.

\keywords{Aprendizagem por Reforço \and Condução Autónoma \and Formula Student \and Soft Actor-Critic \and Simulação}
\end{abstract}

\section{Introdução e Definição do Problema}

A competição Formula Student (FS) é a maior competição de engenharia automóvel universitária a nível mundial. A categoria \textit{Driverless} (FS-AI) desafia equipas a desenvolver veículos totalmente autónomos capazes de navegar numa pista delimitada por cones coloridos --- azuis (esquerda), amarelos (direita) e laranja (partida/chegada) --- sem qualquer intervenção humana~\cite{kabzan2019}. O veículo deve detetar cones, planear trajetórias e executar ações de direção e aceleração em tempo real.

A solução proposta assenta em Aprendizagem por Reforço (RL)~\cite{sutton2018}: treinar um agente que aprenda a conduzir maximizando o progresso e minimizando penalizações. Este relatório documenta a metodologia planeada, incluindo ambiente, algoritmos, ferramentas e formas de validação. O código encontra-se em: \url{https://github.com/pedroreis2468/AR}.

\section{Ambiente de Simulação e Geração de Dados}

O treino será conduzido inteiramente em simulação, com um ambiente desenvolvido de raiz seguindo a API \textit{Gymnasium}.

\textbf{Modelo do veículo.} O veículo é modelado com um modelo cinemático de bicicleta referenciado ao centro de gravidade, calibrado para um FS típico ($\sim$250\,kg, \textit{wheelbase} de 1.55\,m):
$\dot{x} = v \cos(\theta + \beta)$, $\dot{y} = v \sin(\theta + \beta)$,
onde $\beta = \arctan(l_r \tan \delta / L)$ é o ângulo de deslizamento lateral. O modelo inclui limite de aceleração combinada ao atrito disponível ($\mu \cdot g$).

\textbf{Geração procedural e perceção.} As pistas são geradas a partir de pontos de controlo numa elipse perturbada, suavizados com \textit{splines} cúbicas periódicas, randomizando a cada episódio. A perceção é simulada por um sensor de cones (FOV: 180°, alcance: 15\,m), que retorna os 3 cones mais próximos de cada lado com ruído gaussiano ($\sim$5\%).

\textbf{Regras FS-AI.} Cones derrubados: +2\,s por cone (removidos da perceção, sem terminação). Cones laranja: +4\,s. Terminação (DOO) apenas se $\geq$10 cones derrubados, desvio $>$5\,m do \textit{centerline}, ou \textit{off-course} $>$2\,s consecutivos.

\section{Metodologia e Algoritmos}

O problema é formulado como um MDP com $\gamma = 0.99$:
\begin{itemize}[nosep,leftmargin=*]
    \item \textbf{Observação} ($\mathbf{s} \in \mathbb{R}^{20}$): estado do veículo (5\,dims), posições relativas dos cones (12\,dims), distância às fronteiras e erro de \textit{heading} (3\,dims).
    \item \textbf{Ação} ($\mathbf{a} \in [-1, 1]^2$): direção normalizada e aceleração.
    \item \textbf{Recompensa}: $r_t = w_p \cdot \Delta_{\mathrm{prog}} - w_s \cdot |\Delta\delta| - w_l \cdot d_{\mathrm{lat}}^2 - w_t \cdot \Delta t - P_{\mathrm{cones}}$, combinando progresso, suavidade, penalidade lateral, custo temporal e penalizações por cones. Os pesos $w_i$ serão ajustados experimentalmente.
\end{itemize}

O algoritmo principal será o \textbf{Soft Actor-Critic (SAC)}~\cite{haarnoja2018}, um método \textit{off-policy} com maximização de entropia, implementado com \textit{Gaussian actor}, \textit{twin Q-networks} e temperatura $\alpha$ auto-ajustada. Como \textit{baseline}, será utilizado o \textbf{PPO} (\textit{on-policy}) via Stable-Baselines3. Aplicar-se-á \textbf{\textit{domain randomization}}~\cite{ulrich2024} variando massa ($\pm$15\%) e atrito ($\pm$20\%).

\section{Ferramentas e Arquitetura de Software}

O desenvolvimento assenta em Python ($\geq$3.10) com: \textbf{Gymnasium} $\geq$0.29 (API do ambiente), \textbf{PyTorch} $\geq$2.0 (redes \textit{Actor-Critic}), \textbf{Stable-Baselines3} $\geq$2.1 (PPO \textit{baseline}), \textbf{PyGame} $\geq$2.5 (renderização 2D), \textbf{TensorBoard} (monitorização) e \textbf{Git+GitHub} (versionamento e entrega).

\section{Estado Atual da Implementação}

À data de entrega deste relatório, a componente de simulação e o agente SAC encontram-se já implementados e funcionais. Concretamente: o ambiente \texttt{FSRacingEnv} está operacional com modelo cinemático, geração procedural de pistas, sensor de cones com ruído e regras FS-AI completas ($\sim$1500 linhas); o agente SAC custom (PyTorch) está implementado com \textit{Gaussian actor}, \textit{twin Q-networks}, \textit{replay buffer} e auto-tuning de $\alpha$ ($\sim$360 linhas); os \textit{scripts} de treino suportam tanto o modo custom como SB3 (SAC e PPO); e o sistema de avaliação e renderização 2D (PyGame) está funcional. O trabalho restante foca-se na fase experimental: treino extensivo dos agentes, \textit{reward shaping}, experiências de \textit{domain randomization}, análise comparativa de resultados e, como extensão exploratória, investigar a transferência de políticas para o simulador 3D PacSim (\textit{sim-to-sim transfer}).

\section{Protocolo de Validação}

A validação estrutura-se em três vertentes: (1)~\textbf{Desempenho desportivo} --- cada agente avaliado em 20 pistas nunca vistas, reportando taxa de voltas completas, tempo corrigido (bruto + penalizações FS-AI), cones derrubados e velocidade média (média $\pm$ desvio-padrão); (2)~\textbf{Eficiência de aprendizagem} --- monitorização ao longo de 500k \textit{steps} (SAC) e 1M \textit{steps} (PPO) da \textit{reward} cumulativa, \textit{episode length} e \textit{Q-value} estimado; (3)~\textbf{Comparação SAC vs.\ PPO} --- contraste direto em desempenho final, \textit{sample efficiency}, tempo de treino e impacto do \textit{domain randomization} na generalização.

\section{Cronograma}

\begin{table}[t]
\centering
\caption{Planeamento e datas de entrega.}
\label{tab:cronograma}
\begin{tabular}{ll}
\toprule
\textbf{Período} & \textbf{Tarefa} \\
\midrule
Até 30 abr & Ambiente, agente SAC e \textit{pipeline} de treino (\textbf{concluído}) \\
Até 30 abr & \textbf{Entrega do relatório de planeamento} \\
1--11 mai  & Treino intensivo (SAC + PPO) e \textit{reward shaping} \\
12--18 mai & Experiências de \textit{domain randomization} \\
19--25 mai & Análise de resultados; exploração PacSim (sim-to-sim) \\
26 mai--1 jun & Preparação da apresentação \\
2 jun      & \textbf{Apresentação e resultados preliminares} \\
3--15 jun  & Incorporação de feedback; redação do relatório final \\
16 jun     & \textbf{Entrega do relatório final (LNCS)} \\
\bottomrule
\end{tabular}
\end{table}

\begin{thebibliography}{4}

\bibitem{kabzan2019}
Kabzan, J., et al.: AMZ Driverless: The full autonomous racing system. Journal of Field Robotics \textbf{37}(7), 1267--1294 (2019)

\bibitem{sutton2018}
Sutton, R.S., Barto, A.G.: Reinforcement Learning: An Introduction. 2nd edn. MIT Press, Cambridge, MA (2018)

\bibitem{haarnoja2018}
Haarnoja, T., et al.: Soft Actor-Critic algorithms and applications. arXiv preprint arXiv:1812.05905 (2018)

\bibitem{ulrich2024}
Ulrich, F., Wehrli, T.: End-to-End deep reinforcement learning for autonomous racing dynamics. Bachelor Thesis, ZHAW School of Engineering (2024)

\end{thebibliography}

\end{document}
